% ============================================================
% S5-W4 阶段一② 试排样张 —— 卷件「题后紧跟答案」三栏回流制·单页试排
% 2026-09-14｜执行：S5-W4 根因修阶段一臂｜零 git｜写入域＝_tmpS5W4卷件紧跟0914/
% 底本：成卷/测评卷/main.tex 件型层 A~D 照抄（几何/角饰/卷头/页脚零改动）；
%       Q1~Q8 题面与答案逐字迁自同件（原案B 卷末附卷 ansblock 原样内联到题后）。
% 改造点（唯一）＝件型层 E′ 三栏排布器：
%   旧 \jpthree{c1}{c2}{c3}＝\raisebox{0pt}[0pt][0pt]{\vtop{…}}×3 —— 零高固定栏，
%     页构建器视栏高为零 → 答案块入流把栏内容顶过纸边时零报错静默吞（题6/12 实证）。
%   新 \juancols＝multicols 三栏回流 —— 题面与 ansblock 同流顺写，multicol 按
%     真实高度切栏翻页，越限必报 Overfull；栏宽 120／栏距 11／无栏线＝观感参数同旧。
% 编译门：xelatex main.tex ×2 / main-pure.tex ×2（sishou remember picture 硬依赖）。
% ============================================================
\documentclass[fontset=none]{ctexart}
\usepackage[pure]{qp-m3}
\ansblockgrayfalse   % 测评/滚动＝括线模（承旧卷口径；灰底模开关保留不动）
\providecommand{\ov}[1]{\overrightarrow{#1}}

% ------------------------------------------------------------
% 件型层 A｜页面：8 开横放（照抄测评卷，零改动）
% ------------------------------------------------------------
\geometry{paperwidth=399.8mm,paperheight=284.2mm,left=8.9mm,right=8.9mm,
  top=12.0mm,bottom=17.6mm,twoside,headheight=0pt,headsep=0pt,footskip=7.1mm}
\setlength{\topskip}{0pt}
\raggedbottom
\renewcommand{\normalsize}{\fontsize{10.5pt}{17.7pt}\selectfont}
\linespread{1}\normalsize
\setlength{\parindent}{0pt}
\setlength{\parskip}{0pt}
\newlength{\jpcolw}\setlength{\jpcolw}{120.0mm}
\newlength{\jpgap}\setlength{\jpgap}{11.0mm}
\xeCJKsetup{CJKglue={\hskip 0.04em plus 0.04em minus 0.01em}}

\definecolor{silklight}{gray}{0.8235}
\definecolor{silkdark}{gray}{0.6510}
\definecolor{juanrule}{gray}{0.1255}
\definecolor{subgrey}{gray}{0.4}

% ------------------------------------------------------------
% 件型层 B｜JT-01 丝带角饰（照抄测评卷）
% ------------------------------------------------------------
\newcommand{\juanmathSize}{17pt}
\newcommand{\sishou}{%
  \begin{tikzpicture}[overlay,remember picture,x=1mm,y=1mm]
    \path (current page.north west) coordinate (O);
    \begin{scope}[shift={(O)}]
      \fill[silklight] (6.9,-9.39) -- (40.1,-9.39) -- (6.9,-42.59) -- cycle;
      \draw[white,line width=0.28mm] (9.21,-40.29) -- (9.21,-11.63) -- (37.87,-11.63);
      \fill[silkdark] (22.94,-8.21) -- (34.15,-8.21) -- (5.78,-36.58) -- (5.78,-25.37) -- cycle;
      \fill[silkdark] (34.15,-8.21) -- (34.15,-9.20) -- (33.16,-9.20) -- cycle;
      \fill[silkdark] (5.78,-36.58) -- (6.78,-36.58) -- (6.78,-35.58) -- cycle;
      \node[white,inner sep=0pt,anchor=center,rotate=45] at (17.23,-19.74)
        {\fontsize{\juanmathSize}{\juanmathSize}\selectfont\heitiao 数学};
    \end{scope}
  \end{tikzpicture}}

% ------------------------------------------------------------
% 件型层 C｜卷头零件＋题区零件（照抄测评卷）
% ------------------------------------------------------------
\newcommand{\zeroheight}[1]{\raisebox{0pt}[0pt][0pt]{#1}}
\newcommand{\jpcen}[1]{\hbox to\jpcolw{\hfil\zeroheight{#1}\hfil}\nointerlineskip}
\newcommand{\jpright}[1]{\hbox to\jpcolw{\hfill\zeroheight{#1}}\nointerlineskip}
\newcommand{\vgt}[1]{\par\vskip\dimexpr#1-6.25mm\relax}

\newcommand{\tianxieg}[1]{{\fontsize{\qpJTbLabelSize}{13pt}\selectfont\heitiao #1：}%
  \rule[-0.62mm]{\qpJTbLineLen}{0.2mm}}
\newcommand{\jtianxie}{\tianxieg{班级}\hspace{\qpJTbGapGrp}\tianxieg{姓名}%
  \hspace{\qpJTbGapGrp}\tianxieg{得分}}
\newcommand{\juanmingSize}{20.5pt}
\newcommand{\jjuanming}{{\fontsize{\juanmingSize}{26pt}\selectfont\heizhang
  \renewcommand{\CJKglue}{\hskip 0pt}单元素养测评卷（二）}}
\newcommand{\jjunrule}{{\color{juanrule}\rule{\qpJTcRuleLen}{\qpJTcRule}}}
\newcommand{\jfuti}{{\fontsize{\qpJTdSize}{18pt}\selectfont\heitiao\color{subgrey}%
  \renewcommand{\CJKglue}{\hskip 0pt}第二章}}
\newcommand{\jptime}{{\fontsize{\qpJTeSize}{17.7pt}\selectfont（时间：120分钟\quad 分值：150分）}}
\newcommand{\zuhang}[2]{\par\noindent\hangindent=5.7mm\hangafter=1
  {\fontsize{\qpJTfSize}{17.7pt}\selectfont\heitiao #1}#2\par}
\newcommand{\ti}[2]{\par\noindent\hangindent=5.7mm\hangafter=1
  {\fontsize{11.4pt}{17.7pt}\selectfont\heihao{\numboldjian #1}.}\hspace{2.7mm}#2\par}
\newcommand{\fenzhi}[1]{（#1分）}
\newcommand{\optline}[1]{\par\noindent\hangindent=5.7mm\hangafter=0
  {\fontsize{10.5pt}{19pt}\selectfont #1}\par}
\newcommand{\optII}[2]{\makebox[53.05mm][l]{#1}#2}
\newcommand{\optIV}[4]{\makebox[21.5mm][l]{#1}\makebox[21.5mm][l]{#2}\makebox[21.5mm][l]{#3}#4}
\newcommand{\kw}{\hfill（\hspace{3.6mm}）\hspace*{5.4mm}}
\newcommand{\qpart}[1]{\par\noindent\hangindent=5.7mm\hangafter=0
  \hspace*{5.7mm}{\fontsize{10.5pt}{17.7pt}\selectfont #1}\par}

\newdimen\hA \hA=3.03mm
\newdimen\hB \hB=12.12mm
\newdimen\hC \hC=1.89mm
\newdimen\hD \hD=5.95mm
\newdimen\hE \hE=8.04mm
\newdimen\hF \hF=5.70mm
\newdimen\hG \hG=6.35mm
\newcommand{\jphead}{\hbox to\jpcolw{\sishou\hfill}\nointerlineskip
  \vskip\hA\jpright{\jtianxie}%
  \vskip\hB\jpcen{\jjuanming}%
  \vskip\hC\jpcen{\jjunrule}%
  \vskip\hD\jpcen{\jfuti}%
  \vskip\hE\jpcen{\jptime}%
  \vskip\hF
  \zuhang{一、选择题}{：本题共8小题，每小题5分，共40分．\ 在每小题给出的四个选项中，只有一项是符合题目要求的．}%
  \vgt{\hG}}

% ------------------------------------------------------------
% 件型层 D｜YJ-02 文字行页脚（照抄测评卷）
% ------------------------------------------------------------
\newcommand{\jpnum}[1]{{\fontsize{\qpYJbNumSize}{\qpYJbNumSize}\selectfont\numbold #1}}
\newcommand{\jptxtsize}{7.5pt}
\newcommand{\jptxt}{\fontsize{\jptxtsize}{\jptxtsize}\selectfont}
\newcommand{\juanname}{单元素养测评卷（二）}
\newcommand{\pqt}{\ifnum\value{page}<10 0\fi\thepage}
\setlength{\headwidth}{\textwidth}
\fancyfoot[RO]{\jptxt\juanname\hspace{3.95mm}{\heitiao 测评卷}\hspace{3.88mm}卷\hspace{2.25mm}\jpnum{\pqt}}
\fancyfoot[LE]{\jptxt\jpnum{\pqt}\hspace{1.9mm}卷\hspace{3.95mm}{\heitiao 羿郭工作室}%
  \hspace{3.0mm}高中数学\hspace{2.75mm}选择性必修第一册\hspace{2.8mm}RJB}

% ------------------------------------------------------------
% 件型层 E′｜三栏排布器（本试排唯一新件：零高固定栏 → multicols 三栏回流）
%   观感三参数锁定：栏宽 \jpcolw=120mm（＝multicol 列宽）、栏距 \jpgap=11mm（＝\columnsep，
%   同旧手拼 \hspace{\jpgap}）、无栏线（旧手拼三栏本无线）。全幅 382＝3×120＋2×11，
%   三栏逐位等宽等距＝旧观感。multicol 系 qp-m3.sty 既载（\multicolsep=0pt＋\raggedcolumns）。
%   旧 \jpcol 内「\setlength\linewidth{\jpcolw} 换装补钉」退役：multicol 逐栏自更新 \linewidth。
% ------------------------------------------------------------
\setlength{\columnsep}{\jpgap}
\setlength{\columnseprule}{0pt}
\newenvironment{juancols}
  {\begin{multicols}{3}\fontsize{10.5pt}{17.7pt}\selectfont}%
  {\end{multicols}}

\begin{document}
% ============ 题-答-解析同流：题后紧跟 ansblock，三栏自动回流 ============
\begin{juancols}
\jphead
\ti{1}{已知点 \(A(-1,4)\)，\(B(7,-2)\)，则线段 \(AB\) 的中点坐标与 \(|AB|\) 的值分别为\kw}
\optline{\optII{A．\((3,1)\)，\(10\)；}{B．\((6,2)\)，\(10\)；}}
\optline{\optII{C．\((3,1)\)，\(100\)；}{D．\((4,2)\)，\(10\)．}}
\begin{ansblock}[2章-测-1]
% ans:2章-测-1
\ansitem{1}{A}
\ansnote{解析}{中点 \(M=\left(\dfrac{-1+7}{2},\dfrac{4-2}{2}\right)=(3,1)\)；\(|AB|=\sqrt{8^{2}+(-6)^{2}}=10\)。B 系中点漏除以 \(2\)；C 系忘开方；D 系中点算误。故选 A。}
\end{ansblock}
\ti{2}{下列说法中正确的是\kw}
\optline{A．若两条直线斜率相等，则它们互相平行；}
\optline{B．若 \(l_{1}\parallel l_{2}\)，则 \(k_{l_{1}}=k_{l_{2}}\)；}
\optline{C．若两条直线中有一条直线的斜率不存在，另一条直线的斜率存在，则这两条直线相交；}
\optline{D．若两条直线的斜率都不存在，则它们相互平行．}
\begin{ansblock}[2章-测-2]
% ans:2章-测-2
\ansitem{2}{C}
\end{ansblock}
\ti{3}{已知点 \(A(1,4)\)，\(B(5,2)\)，点 \(P\) 为直线 \(l\colon x+y=0\) 上的动点，则 \(|PA|+|PB|\) 的最小值为\kw}
\optline{\optIV{A．\(3\sqrt{10}\)；}{B．\(2\sqrt{5}\)；}{C．\(\sqrt{34}\)；}{D．\(3\sqrt{5}\)．}}
\begin{ansblock}[2章-测-3]
% ans:2章-测-3
\ansitem{3}{A}
\ansnote{解析}{\(A\)、\(B\) 代入 \(x+y=0\) 左侧均大于 \(0\)，同侧。设 \(A\) 关于 \(l\) 的对称点 \(A'(-4,-1)\)，则 \(|PA|+|PB|=|PA'|+|PB|\geq|A'B|\)，取等点 \(P_{0}\left(-\dfrac{1}{4},\dfrac{1}{4}\right)\in l\)。最小值 \(=|A'B|=\sqrt{9^{2}+3^{2}}=3\sqrt{10}\)。故选 A。}
\end{ansblock}
\ti{4}{与 \(y\) 轴相交于 \(A(0,2)\)、\(B(0,6)\) 两点，且半径等于 \(2\sqrt{2}\) 的圆的方程是\kw}
\optline{\optII{A．\((x+2)^{2}+(y-4)^{2}=8\)；}{B．\((x-2)^{2}+(y-4)^{2}=8\)；}}
\optline{C．\((x+2)^{2}+(y-4)^{2}=8\) 或 \((x-2)^{2}+(y+4)^{2}=8\)；}
\optline{D．\((x-2)^{2}+(y+4)^{2}=8\)．}
\begin{ansblock}[2章-测-4]
% ans:2章-测-4
\ansitem{4}{C（两圆并取）}
\end{ansblock}
\ti{5}{已知两点 \(A(3,0)\)，\(B(0,3)\)，动点 \(P\) 满足 \(|PA|^{2}+|PB|^{2}=10\)，则 \(P\) 的轨迹方程为\kw}
\optline{\optII{A．\(x^{2}+y^{2}-3x-3y+4=0\)；}{B．\(x^{2}+y^{2}-3x-3y-4=0\)；}}
\optline{\optII{C．\(x^{2}+y^{2}+3x+3y+4=0\)；}{D．\(x^{2}+y^{2}-3x-3y+10=0\)．}}
\begin{ansblock}[2章-测-5]
% ans:2章-测-5
\ansitem{5}{A}
\ansnote{解析}{设 \(P(x,y)\)：\((x-3)^{2}+y^{2}+x^{2}+(y-3)^{2}=10\)，化简得 \(x^{2}+y^{2}-3x-3y+4=0\)，即 \(\left(x-\dfrac{3}{2}\right)^{2}+\left(y-\dfrac{3}{2}\right)^{2}=\dfrac{1}{2}\)。全过程同解变形，且 \(A\)、\(B\) 均不满足条件，纯粹性与完备性同时成立。故选 A。}
\end{ansblock}
\ti{6}{到两定点 \(F_{1}(-3,0)\)，\(F_{2}(3,0)\) 的距离之差的绝对值等于 \(6\) 的点 \(M\) 的轨迹为\kw}
\optline{\optIV{A．椭圆；}{B．两条射线；}{C．双曲线；}{D．线段．}}
\begin{ansblock}[2章-测-6]
% ans:2章-测-6
\ansitem{6}{B（两条射线）}
\end{ansblock}
\ti{7}{若抛物线 \(y^{2}=\dfrac{4}{m}x\) 的焦点与椭圆 \(\dfrac{x^{2}}{7}+\dfrac{y^{2}}{3}=1\) 的左焦点重合，则 \(m\) 的值为\kw}
\optline{\optIV{A．\(-\dfrac{1}{2}\)；}{B．\(\dfrac{1}{2}\)；}{C．\(-2\)；}{D．\(2\)．}}
\begin{ansblock}[2章-测-7]
% ans:2章-测-7
\ansitem{7}{A（\(m=-\dfrac{1}{2}\)）}
\end{ansblock}
\ti{8}{已知 \(F_{1}\)，\(F_{2}\) 分别是双曲线 \(C\colon \dfrac{x^{2}}{3}-\dfrac{y^{2}}{9}=1\) 的左、右焦点，过 \(F_{2}\) 的直线与双曲线 \(C\) 的右支交于 \(A\)，\(B\) 两点，\(\triangle AF_{1}F_{2}\) 和 \(\triangle BF_{1}F_{2}\) 的内心分别为 \(M\)，\(N\)，则 \(|MN|\) 的取值范围是\kw}
\optline{\optIV{A．\([2\sqrt{3},4)\)；}{B．\([2\sqrt{3},3)\)；}{C．\([2\sqrt{3},+\infty)\)；}{D．\((2\sqrt{3},3)\)．}}
\begin{ansblock}[2章-测-8]
% ans:2章-测-8
\ansitem{8}{A（\([2\sqrt{3},4)\)）}
\end{ansblock}
\end{juancols}
\end{document}
